\begin{pastexam}[label={px:ccu-114-csb-05}]{中正資工所乙組 114 學年度 第 5 題}
（是非組合；單選）

(5\%) Please answer the following \textbf{true-or-false} questions and select the correct set of answers from the provided options.
\begin{enumerate}[label=\Roman*.]
\item The adjacency matrix representation of a graph is more memory-efficient than the adjacency list for sparse graphs.
\item Breadth-first search (BFS) always finds the shortest path in weighted graphs.
\item A new node that is inserted into a binary search tree is not always at a leaf node.
\end{enumerate}

A.~False, False, False \qquad B.~False, False, True \qquad C.~False, True, False\\
D.~False, True, True \qquad E.~True, False, False \qquad F.~True, False, True\\
G.~True, Ture, False \qquad H.~Ture, True, True
\end{pastexam}

\begin{pastsoln}
答案是 A（False, False, False）。 I、II 是資結書第~21~章、III 是第~9~章的內容；歸在第~21~章。

\begin{itemize}
\item \textbf{I：False。} 相鄰矩陣不管邊有多少，都用 $\Theta(n^2)$ 的空間；相鄰串列用 $\Theta(n + m)$（資結書第~21.2~節）。 稀疏圖
$m$ 遠小於 $n^2$，例如 $m = O(n)$，相鄰串列只要 $\Theta(n)$，比矩陣省得多。 說反了。
\item \textbf{II：False。} BFS 找到的是\textbf{邊數}最少的路徑（資結書第~21.3~節），只有在每條邊權重都相同時才是最短路徑。 反例：
$s \to t$ 權重 $10$，$s \to x \to t$ 權重 $1 + 1$。 BFS 從 $s$ 出發一步就到 $t$，給出長 $10$ 的路徑，最短的是 $2$。 有權重時要用
Dijkstra 之類的演算法。
\item \textbf{III：False。} 二元搜尋樹的插入從根往下找，走到空的位置才把新節點接上（資結書第~9.2~節），新節點一定是葉節點。
敘述說「不一定在葉節點」，錯。（新節點插入後可能讓原本的葉節點變成內部節點，但新節點本身插入當下一定是葉。）
\end{itemize}
三個都是 False，選 A。
\end{pastsoln}
